\subsection{避孕责任：谁承担、怎么谈}

避孕的技术问题有清晰答案，责任问题没有。前面几节讲了方法、效果与差错处理，但还有一个几乎从不被写进任何说明书的问题：\textbf{这件事该由谁负责？}本节处理三个具体问题——现在实际是谁在承担、伴侣间怎么谈、谈不拢时怎么办。

\textbf{一、现状：避孕是一项被默认分派给女性的工作}

\begin{itemize}
\item 在可用的避孕方法中，绝大多数作用于女性身体：口服药、贴片、阴道环、注射剂、皮下埋植、宫内节育器、女性绝育手术。男性可用的常规方法主要只有安全套与绝育手术，体外射精则效率很低；
\item 因此副作用也主要由女性承担：激素相关的不适与情绪波动、放置或取出节育器时的疼痛、不规则出血、体重变化，以及绝育手术在侵入性上明显的性别差异；
\item 这种分工常被当成"自然如此"，但它更多是技术发展史与社会默认值的结果，而不是生物学上的必然。把它看成"默认值"而不是"义务"，是讨论这件事的起点。
\end{itemize}

\textbf{二、男性可以实际承担的部分（清单）}

这一节的立场不是"把责任互换"，而是"把责任分担"：

\begin{itemize}
\item \textbf{使用并自行准备安全套}：包括自己购买、随身携带、了解尺寸与正确用法（见本书避孕套相关小节），而不是"临到用的时候问对方有没有"；
\item \textbf{共同承担"记得"这件事}：帮助记录、提醒、取药、陪同就诊，把"每天都得想着"的认知负担从一个人身上分出来一部分；
\item \textbf{在伴侣出现副作用时参与决策}：当女方因激素不适、出血、疼痛而考虑更换方法时，参与讨论与陪同就诊，而不是把这件事留给她自己处理；
\item \textbf{不做甩手掌柜式的知情}：了解伴侣所使用的方法、它的失败率与需要加用屏障的情形，这在漏服、腹泻、呕吐等情况下会直接决定是否需要补救措施；
\item \textbf{不把欲望下降归咎于对方}：部分女性在服用激素避孕后性欲下降（机制见女性卷相关小节），这是药物效应，不是"不爱了"；
\item \textbf{长期规划上保持对称}：如果双方已确定不再生育，绝育的选项应当在两人之间被讨论，而不是默认由女方承担手术。
\end{itemize}

\textbf{三、当伴侣不愿意使用安全套：可用的回应}

这是实际问题中最常出现、也最容易被道德指责替代的一段。比较有用的做法是把三件事分开：

\begin{itemize}
\item \textbf{把"不舒服"变成可解决的技术问题}：合适尺寸（过紧会勒痛、过松易滑脱）、超薄型、配合水性润滑剂，都能明显改善体验。多数"戴套没感觉"的抱怨，实际来源是尺寸不合与润滑不足，而不是安全套本身；
\item \textbf{把选项摊开，而不是只剩一个}：如果对方抗拒安全套，可以讨论其他共同承担的方式，例如由女方使用长效方法、或双方共同确认检测状况后采用"一种高效方法 + 定期检测"的组合。前提是\textbf{这些替代方案是双方共同接受的，不是被说服接受的}；
\item \textbf{明确边界，并知道这句话不需要理由}："我不在没有安全套的情况下进行"——这是一句完整的话，不需要补充解释。拒绝不构成对关系的伤害，也不需要"给个说法"；
\item \textbf{不要接受的两种情形}：一是以"你不信任我"施压，把使用安全套等同于怀疑感情；二是\textbf{在对方不知情或不同意的情况下取下安全套}（部分文献称 stealthing）。后者不是"小动作"或"情趣"，而是对同意范围的直接侵害——同意的内容是带安全套的行为，擅自改变前提条件，等同于在原同意范围之外行事（参见本书性同意与性侵犯相关章节）。
\end{itemize}

\textbf{四、把避孕从"谈判"变成"共同管理"}

\begin{itemize}
\item \textbf{设定复盘节奏}：每 6--12 个月一起看一次——当前方法是否合适、副作用是否可接受、生育计划是否有变化。避孕方法的选择本来就不是一次性的决定；
\item \textbf{共同记录关键信息}：方法名称、开始时间、需要更换或取出的时间点（如埋植剂与节育器的有效期、注射剂的下一次时间）；
\item \textbf{计划怀孕时一起准备}：提前停药或取出、调整叶酸等，都是两个人的事；
\item \textbf{把紧急情况写成预案}：如果发生避孕套破裂或漏服，第一步做什么（见本书相关小节）——事先说一次，比事后争执有用得多。
\end{itemize}

\textit{【编者注】}本节不对关系形式（婚前、婚内、非单一伴侣）作预设，只讨论一个技术性的公平问题：避孕的风险由谁的身体承担，决定就应当由谁参与。临床上有一个值得记住的现象——当避孕由一方独自管理时，"忘记""不敢说""将就一次"的概率都会上升；而当它成为两人共同的事务时，正确使用的比例会明显提高。

